\subsection{收缩描述}

第四问根据给定半径轨迹重新确定烘干时长，并更换另一组经验物性。半径记录覆盖 $0$--$72\,\mathrm{h}$，由 $2\,\mathrm{cm}$ 单调降至 $1.198\,\mathrm{cm}$。附件只给出外半径，不能唯一识别内部非均匀变形。本文采用均匀比例径向收缩，并取长度 $L$ 不变。半径函数 $R(t)$ 对观测点作 PCHIP，保证 $R>0$ 且不产生非物理反弹。

干基含水率 $C$ 是水与干物质的质量比，不是单位体积含水量。在固定长度、干物质不流失且初始干物质密度均匀的条件下，径向固体速度取
\[
v(r,t)=\frac{\dot R(t)}{R(t)}r.
\]
引入材料坐标 $\xi=r/R(t)$。均匀收缩使材料速度与网格速度一致，对流项抵消，水分方程化为固定单位区间上的变系数扩散：
\begin{equation}
\frac{\partial U}{\partial t}
=\frac{1}{R(t)^2\xi}\frac{\partial}{\partial\xi}
\left[\xi D(U,\Theta)\frac{\partial U}{\partial\xi}\right],
\qquad 0<\xi<1,
\label{eq:q4-C}
\end{equation}
其中 $U(\xi,t)=C(R(t)\xi,t)$。相应的有效能量方程为
\begin{equation}
\rho(U)c_p(U)\frac{\partial\Theta}{\partial t}
=\frac{1}{R(t)^2\xi}\frac{\partial}{\partial\xi}
\left[\xi k(U)\frac{\partial\Theta}{\partial\xi}\right].
\label{eq:q4-T}
\end{equation}
不能只移动边界而漏掉材料运动，也不能在已经抵消后再重复加入几何对流项\cite{adrover2019}。

\subsection{物性与边界}

附录给出
\begin{equation}
\begin{aligned}
\rho&=760+90U,\\
c_p&=1850+2150\frac{U}{1+U},\\
k&=0.12+0.20\frac{U}{1+U},\\
D&=4.2\times10^{-4}\exp(-0.30/U)\exp\bigl[-3850/(\Theta+273.15)\bigr].
\end{aligned}
\label{eq:q4-prop}
\end{equation}
初值 $U=2.55$、$\Theta=28^\circ\mathrm{C}$，中心导数为零。表面等效边界写到 $\xi=1$ 上：
\[
-\frac{k}{R}\frac{\partial\Theta}{\partial\xi}=h(\Theta_s-T_a),\qquad
-\frac{D}{R}\frac{\partial U}{\partial\xi}=h_m(U_s-C_a).
\]
环境与问题三相同：$0$--$4\,\mathrm{h}$ 用观测曲线，之后取 $50^\circ\mathrm{C}$、$0.05\,\mathrm{kg/kg}$。$\rho c_p$ 仍解释为有效热容量；若把同一 $\rho$ 同时当成真实湿密度，并与固定长度、均匀收缩和给定 $R(t)$ 联立，会在后期过约束，因此不以该解释作为严格多相守恒模型。

同一温度、含水率下，本问扩散系数与问题三所用公式之比约为 $0.175\exp(0.15/C)$，在湿态约为 $0.19$、在 $C=0.15$ 时约为 $0.48$。收缩使扩散距离缩短，但物性同时变慢，第三、四问时长之差不能全部归因于半径变化。

\subsection{达标与输出}

达标判据与\eqref{eq:q3-event}相同，最大值在 $\xi\in[0,1]$ 上取。输出按当前物理半径：$r=\xi R(t)$。表~\ref{tab:q4C} 在每 $6\,\mathrm{h}$ 及结束时刻给出 $0,0.5,\ldots\,\mathrm{cm}$ 且不超过当时半径的位置，并单列实际表面值；域外单元格留空，不得填 $0$ 或空气含湿量。例如 $6\,\mathrm{h}$ 时 $R=1.374\,\mathrm{cm}$，有效列仅为 $0$、$0.5$、$1.0\,\mathrm{cm}$ 及表面。

求解仍在固定 $\xi$ 网格上作有限体积与联合 BDF，每步更新 $R(t)$、$k$、$D$ 和 $\rho c_p$。退化检验包括：$R$ 为常数时回到固定域模型；关闭边界交换且 $C$ 均匀时，纯收缩不改变干基含水率。结束时刻的数值同样待积分完成后填入。

\begin{table}[htbp]
\centering
\caption{收缩条件下烘干过程的水分浓度（单位：$\mathrm{kg/kg}$，结束时刻待补）}
\label{tab:q4C}
\begin{tabular}{cccccc}
\toprule
时间/h & $0\,\mathrm{cm}$ & $0.5\,\mathrm{cm}$ & $1.0\,\mathrm{cm}$ & $\cdots$ & 药材表面 \\
\midrule
6  & --- & --- & --- & --- & --- \\
12 & --- & --- & --- & --- & --- \\
18 & --- & --- & --- & --- & --- \\
$\cdots$ & --- & --- & --- & --- & --- \\
结束时刻 & --- & --- & --- & --- & --- \\
\bottomrule
\end{tabular}
\end{table}
